\chapter*{Введение}\label{body-start}
\addcontentsline{toc}{chapter}{Введение}

Автоматические тесты регулярно запускаются во время разработки. Когда проверок становится много, время ожидания, загрузка процессора и потребление памяти начинают влиять на удобство работы и требования к машине. Поэтому выбор тестового фреймворка полезно обосновывать результатами на собственном приложении.

В работе сравниваются Jest, Vitest и Mocha. Цель -- выбрать наиболее подходящий фреймворк для тестирования разработанного консольного приложения по времени выполнения полного набора тестов, накопленному CPU-времени и наблюдаемому пику оперативной памяти. Все три инструмента получают одинаковые входные данные и выполняют одинаковые проверки.

Приложение написано на JavaScript и реализует алгоритм Манакера. Оно принимает строку из консоли или список строк из JSON-файла. Для каждого входа выводятся максимальный палиндром, его длина и число палиндромных вхождений. Выбран один простой формат файла, чтобы код импорта и правила проверки были понятны при чтении.

Основное внимание уделяется тестированию. Подготовлены 99 отдельных тестов для каждого фреймворка: примеры с явными ожиданиями, граничные значения, независимый перебор коротких строк, свойства результата, ошибки JSON, файловый ввод и реальные запуски консоли. Дополнительно в изолированные копии программы вносятся контрольные дефекты, чтобы проверить чувствительность тестового набора.

Для сравнительного эксперимента разработан внешний наблюдатель ресурсов. Он последовательно запускает раннеры, сохраняет журналы, проверяет число пройденных тестов и записывает исходные результаты. Размер JSON-файла изменяется за счёт длины содержащихся в нём строк. На каждом размере выполняются повторные измерения, из которых рассчитываются медианы и квартили.

Первая глава описывает проблему, критерии сравнения, алгоритм и порядок эксперимента. Вторая посвящена приложению, структуре тестов и командам запуска. Третья содержит условия апробации и измеренные показатели. В четвёртой сопоставляются результаты и обосновывается выбор фреймворка с учётом границ проведённого эксперимента.
